%-----------------------------------------------------------------------
\section{Discussion}
%-----------------------------------------------------------------------

The headline 69.04\% top-1 accuracy is encouraging, but accuracy is not the
result that matters most here. The decisive change is sequential: earlier
versions of the pipeline labelled over half of all windows \texttt{none},
shredding the shot-to-shot grammar, whereas improved pose features and
classifier fine-tuning cut the \texttt{none} rate to 6.70\% and make the rally
continuous again. Only on a continuous stream do the mining views become
legible and---per the validation protocol
(Table~\ref{tab:mining_eval})---demonstrably faithful rather than merely
plausible. We therefore spend this
section reading the recovered grammar the way a coach would: granting that the
numbers are trustworthy, what do they actually say about how elite singles
points are built?

\textbf{The point is built, not struck---the smash is pressure, not a finisher.}
The most counter-intuitive finding is that at this level the smash almost never
ends the rally: 61.9\% of smashes come straight back as a defensive return net,
with fewer than 3\% drawing a lob or net-shot reply---the rally continues in
almost every case. Elite
defence is good enough that one smash rarely beats it, so the smash functions to
\emph{extract} a weak reply that sets up the next stroke rather than to finish
outright. The natural tactical unit is thus the sequence
smash~$\rightarrow$~block~$\rightarrow$~kill, not the smash in isolation---which
is also why a method that preserves sequence coherence, not just per-stroke
accuracy, is the right instrument for studying the sport. For practitioners this
argues for drilling the smash-and-follow-up and valuing recovery speed after
smashing as highly as raw smash power.

\textbf{Attacks are manufactured at the net, not generated by power.}
The grammar locates the origin of attack in the forecourt. Net shots force a
lift 52.0\% of the time (and a push 20.9\%), and a lift is then answered with a
smash or drop 70.7\% of the time; drops likewise route about half of all replies
back into the net zone. The chain that ends in a smash therefore usually
\emph{begins} with a tight net shot or drop that compels the opponent to lift.
Power closes the phrase, but net skill opens it, so winning the net exchange is
the precondition for ever getting to attack.

\textbf{The lift is a concession: net play is a contest over who lifts first.}
Because lifting hands initiative away (lob~$\rightarrow$~smash 41.6\%,
lob~$\rightarrow$~drop 29.1\%), the forecourt exchange is effectively a duel to
\emph{avoid} lifting: 20.3\% of net shots draw another net shot, sustaining a
net-zone standoff in which each player tries to force the other up first. This
recasts the net not as a place to score but as a place not to concede, and
explains the premium elite players place on tight, spinning net shots.

\textbf{Elite singles runs in two gears---aggressive engagement and a patient
waiting game.} The transitions expose two stable cycles rather than one. The
engagement loop net~$\rightarrow$~lob~$\rightarrow$~smash~$\rightarrow$~return
net~$\rightarrow$~push/lob is the attacking rhythm; the clear~$\rightarrow$~clear
continuation (30.6\%) is a backcourt holding pattern in which neither player
commits until a weak ball arrives. The game is an alternation between escalation
and patience, and the discipline to keep a rally neutral---rather than forcing a
low-percentage attack from a balanced position---emerges here as a measurable
trait rather than a vague virtue.

\textbf{Mobility is a prerequisite for aggression, not a separate trait.}
The one movement correlation that survives strict multiplicity control links
lateral coverage width to net-ward pressure ($r=+0.73$, $q=0.04$): the
player-sets that cover the most sideways ground are also the ones that press the
net most. Attacking forecourt play appears to be \emph{enabled by} footwork and
court coverage---one can only commit forward while trusting the sideways
recovery---so conditioning and tactics are entangled rather than independent.

\textbf{Most elite player-sets are all-round aggressors; sustained defence is
rare and looks imposed.} Clustering separates an Active-Attacker majority
($8/16$, highest smash and net-attack rates) from a small Passive-Baseliner
group ($3/16$) and a Rear-court-Dominator group ($5/16$) that retrieves deep and
then presses forward. Tellingly, the Passive Baseliners sit \emph{closest} to
the net on average yet attack it \emph{least}: their forward position is pushed
onto them by the opponent rather than chosen, suggesting that at this level
prolonged passivity is something that happens \emph{to} a player, not a strategy
a player selects. We read these archetypes as descriptive (assignment ARI
$0.17$) rather than a transferable taxonomy, but the feature axes beneath them
are faithful ($\bar{r}=0.79$).

\textbf{The far-side bottleneck.}
Feature quality explains most of the residual error. Far-side pose is valid in
only 54.29\% of windows against 98.96\% near side; windows above 75\% far-side
validity reach 76.71\% accuracy versus 56.05\% below 50\%---a gap of more than 20
percentage points that pinpoints far-side extraction as the highest-value target
for future work. Reassuringly, the tactical conclusions above survive this
noise: the predicted-vs-true gap is smallest exactly at the two- and
three-stroke phrase scale that carries the most coaching value, so the grammar
we read off the predicted stream is the part the pipeline recovers best.
